% Gerado por scripts/migrate_markdown_to_latex.py.
% Fonte migrada: docs/relatorio_execucao_autonoma_2026-08-25.md


\chapter{Relatório da execução autônoma - 25/08/2026}



\section{Resumo executivo}







A validação numérica da atenção está concluída, a metodologia pré-benchmark foi auditada e o pipeline de operações isoladas foi implementado, testado e usado em três coletas CPU versionadas. Todos os portões controláveis pelo repositório estão fechados. A coleta principal está corretamente bloqueada porque esta máquina não possui CUDA nem a RTX 4070 Ti Super prevista no plano.







O diagnóstico maior revelou uma ameaça à validade - pesos \texttt{N(0,1)} faziam a escala da saída crescer com \texttt{D}. O resultado original foi preservado, o gerador foi corrigido para Xavier/bias zero e a coleta foi repetida. A hipótese H1 (quantização preserva melhor a saída que pruning de 50\%) foi compatível nos casos executados, mas ainda é preliminar e não equivale a ganho de hardware.



\section{O que foi feito}


\begin{enumerate}
\item Auditoria do repositório, cinco commits preexistentes, código ativo/legado,
documentação, plano formal, apresentação, reunião e ambiente.
\item Arquivamento da transcrição integral e reconstrução classificada da
orientação, separando pedidos, decisões, dúvidas e interpretações.
\item Validação NumPy/PyTorch/TensorFlow-Keras do núcleo da atenção: 9/9 aprovada.
\item Criação do diário, checklist pré-benchmark e diagnóstico metodológico.
\item Correção da semântica de hardware do runner piloto.
\item Novo runner de projeção densa/self-attention com configuração versionada,
lotes independentes, logs, hashes, manifesto e recusa de GPU incorreta.
\item Análise com checksums, baseline por execução, média, desvio-padrão, relatório
e gráficos.
\item Smoke CPU ponta a ponta.
\item Diagnóstico CPU \texttt{L=64}, \texttt{D=128} v1, investigação da escala dos pesos,
correção metodológica e repetição v2.
\item Suíte final observada: \texttt{45 passed, 1 skipped}; \texttt{pip check} sem problemas.
\end{enumerate}


\section{Commits criados nesta execução}



\begin{longtable}{@{}>{\RaggedRight\arraybackslash}p{0.410\linewidth}>{\RaggedRight\arraybackslash}p{0.410\linewidth}@{}}
\toprule
Commit & Etapa lógica \\
\midrule
\endfirsthead
\toprule
Commit & Etapa lógica \\
\midrule
\endhead
\texttt{6e10da0} & Reunião, transcrição e direcionamentos do professor. \\
\texttt{2a42a5e} & Validação da atenção entre PyTorch e TensorFlow/Keras. \\
\texttt{49f081f} & Auditoria, diário e portões pré-benchmark. \\
\texttt{cedbeb2} & Correção das alegações de hardware do piloto. \\
\texttt{7113fa5} & Runner reproduzível das operações isoladas. \\
\texttt{bcea919} & Captura do ambiente antes de escrever artefatos. \\
\texttt{a0aebe5} & Análise entre execuções e validação de integridade. \\
\texttt{1474fb5} & Smoke CPU versionado. \\
\texttt{df061d8} & Configuração do diagnóstico CPU representativo. \\
\texttt{cc03231} & Diagnóstico v1 e registro do risco de inicialização. \\
\texttt{bf7b4bd} & Pesos Xavier e schema de configuração v2. \\
\texttt{4cd1ebf} & Diagnóstico v2 e confronto metodológico v1/v2. \\
\bottomrule
\end{longtable}




Este relatório constitui a etapa final de handoff e será versionado em commit separado após a verificação final.



\section{Orientações recuperadas do professor}



Pedidos explícitos:


\begin{itemize}
\item validar a implementação contra TensorFlow/Keras usando a mesma entrada e
semente antes de otimizar;
\item comparar equação/manual, PyTorch e TensorFlow em tabela ou gráfico;
\item só depois escolher uma ou duas técnicas e comparar os resultados;
\item concentrar o trabalho na self-attention e nas matrizes de maior custo, não em
uma Transformer completa;
\item escrever texto corrido e fechar avaliações com resultados/tabelas desde já.
\end{itemize}


Decisões ou sugestões condicionais:


\begin{itemize}
\item uma operação interna equivalente basta se a camada completa ocultar detalhes;
\item inferência com modelo do Hugging Face é extensão dependente de tempo;
\item novembro foi citado de forma incerta e não substitui a data institucional.
\end{itemize}


Interpretações posteriores, não atribuídas ao professor:


\begin{itemize}
\item manter projeção densa junto da self-attention por constar no plano formal;
\item usar NumPy como referência matemática;
\item escolher Keras Ops, tolerâncias, schema de artefatos e níveis de validade;
\item usar Xavier/bias zero para controlar a escala entre dimensões.
\end{itemize}


\section{Pré-benchmark}



\begin{longtable}{@{}>{\RaggedRight\arraybackslash}p{0.273\linewidth}>{\RaggedRight\arraybackslash}p{0.273\linewidth}>{\RaggedRight\arraybackslash}p{0.273\linewidth}@{}}
\toprule
Portão & Status & Evidência \\
\midrule
\endfirsthead
\toprule
Portão & Status & Evidência \\
\midrule
\endhead
Reunião reconstruída sem inventar requisitos & Concluído & \texttt{docs/\allowbreak{}reunioes/\allowbreak{}2026-07-31/\allowbreak{}}. \\
Atenção validada entre referências & Concluído & 9/9 em \texttt{evidencias/\allowbreak{}comparacao\_\allowbreak{}frameworks/\allowbreak{}}. \\
Recorte e cenários fixados & Concluído & Projeção densa, self-attention, baseline/\allowbreak{}pruning/\allowbreak{}quantização. \\
Configuração formal codificada & Concluído & Seed 42, grade, 20/50, duas execuções. \\
Rastreabilidade e integridade & Concluído & Config, ambiente, log, manifesto, SHA-256 e fingerprints. \\
Análise entre execuções & Concluído & Baseline da mesma execução, média e DP. \\
Inicialização controlada & Concluído & Schema v2, Xavier e confronto v1/v2. \\
Smoke ponta a ponta & Concluído & 12 registros e reprodutibilidade aprovada. \\
GPU-alvo disponível & Bloqueado externamente & Intel Iris Xe; \texttt{nvidia-smi} ausente; CUDA count 0. \\
\bottomrule
\end{longtable}



\section{Benchmarks e validações executados}



\subsection{Comparação da atenção entre frameworks}


\begin{itemize}
\item Objetivo: provar correção do núcleo antes de otimizar.
\item Configuração: três casos, CPU/float32, seed 2026, \texttt{atol=1e-6}, \texttt{rtol=1e-5}.
\item Resultado: 9/9; erro absoluto máximo \texttt{2.\allowbreak{}38418579e-07}.
\item Interpretação: equivalência numérica nos casos testados; nenhuma comparação de
desempenho.
\end{itemize}


\subsection{Smoke CPU}


\begin{itemize}
\item Objetivo: validar pipeline completa.
\item Configuração: \texttt{B=1}, \texttt{L=4}, \texttt{D=8}, 2 warm-ups, 5 medições, 2 execuções.
\item Resultado: 12 registros, checksums e hashes aprovados; gráficos legíveis.
\item Interpretação: infraestrutura aprovada; tempos muito variáveis e sem valor de
hardware. Quantização teve MSE menor que pruning nas duas operações.
\end{itemize}


\subsection{Diagnóstico CPU v1}


\begin{itemize}
\item Objetivo: executar um ponto formal com 20/50 e duas execuções.
\item Configuração: \texttt{B=1}, \texttt{L=64}, \texttt{D=128}, pesos \texttt{N(0,1)}.
\item Resultado: H1 compatível, mas MSE da atenção quantizada \texttt{732.\allowbreak{}65}.
\item Interpretação: resultado legítimo que revelou escala inadequada do baseline;
não foi descartado nem 'melhorado'.
\end{itemize}


\subsection{Diagnóstico CPU v2}


\begin{itemize}
\item Objetivo: repetir o mesmo ponto após controlar a inicialização.
\item Configuração: mesmos parâmetros, pesos Xavier e bias zero.
\item Resultado: na projeção, MSE \texttt{7.\allowbreak{}7789e-05} (quantização) contra \texttt{0.\allowbreak{}0727207}
(pruning); na atenção, \texttt{1.\allowbreak{}63925e-05} contra \texttt{0.\allowbreak{}0110127}. R²/cosseno da atenção
quantizada: \texttt{0.\allowbreak{}999647}/\texttt{0.\allowbreak{}999824}.
\item Interpretação: H1 compatível nesse ponto; escala absoluta controlada. Latência
CPU abaixo do baseline em algumas barras não autoriza aceleração devido ao
ruído e à ausência de kernels especializados.
\end{itemize}


\section{Problemas e riscos encontrados}


\begin{itemize}
\item GPU-alvo indisponível nesta máquina.
\item Runner anterior media bloco inteiro, divergindo do recorte de operações.
\item Defaults anteriores divergiam da seed/grade/20/50/duas execuções do plano.
\item Quantização manual executa pesos dequantizados \texttt{float32}; pruning permanece
denso. Nenhum dos dois sustenta hoje alegação de kernel INT8/esparso.
\item Memória em CPU é estimativa de armazenamento dos tensores, não pico real.
\item Latências pequenas apresentaram alta variabilidade; duas execuções são o
mínimo formal, não uma amostra estatística ampla.
\item Requisitos comuns têm versões mínimas; cada coleta registra o snapshot exato.
\item O plano declara maio-setembro de 2026, mas o cronograma contém uma célula com
'2º semestre de 2027'; a reunião menciona novembro sem confirmação.
\item O \texttt{.\allowbreak{}docx} formal foi extraído estruturalmente, mas não renderizado visualmente
porque LibreOffice não está disponível.
\end{itemize}


\section{Pendências}



\subsection{Críticas}


\begin{itemize}
\item executar a grade principal na RTX 4070 Ti Super com CUDA;
\item antes de alegar ganho físico, adotar e validar formato/kernel realmente INT8
e/ou esparso; os cenários simulados atuais só sustentam validade numérica ou
algorítmica.
\end{itemize}


\subsection{Importantes}


\begin{itemize}
\item confirmar o calendário institucional e a inconsistência 2026/2027;
\item incorporar tabelas e discussão v1/v2 ao texto acadêmico;
\item após a coleta GPU, revisar ruído, outliers e suficiência de duas execuções;
\item decidir, com o orientador, se a entrega formal manterá quantização/pruning
simulados ou incluirá uma extensão de kernel real.
\end{itemize}


\subsection{Melhorias futuras}


\begin{itemize}
\item modelo pré-treinado do Hugging Face;
\item TorchAO/kernel INT8 e esparsidade estruturada em etapa própria;
\item mais seeds/lotes/heads e intervalos de confiança;
\item benchmark do bloco completo somente como extensão;
\item comparação de desempenho entre frameworks, separada da validação numérica.
\end{itemize}


\section{Próximo passo recomendado}







Mover o repositório para a máquina com RTX 4070 Ti Super, instalar o mesmo ambiente, rodar primeiro a suíte e o smoke v2, confirmar o nome da GPU e então executar \texttt{experimentos/\allowbreak{}benchmark\_\allowbreak{}principal\_\allowbreak{}gpu.\allowbreak{}json} em pasta nova. A coleta deve ser analisada imediatamente e não pode ser chamada de ganho INT8/esparso sem uma segunda etapa que comprove armazenamento e kernel compatíveis.
